\documentclass[11pt]{article}
\usepackage{hwtools}

\newcommand{\C}{\mathbb{C}}
\newcommand{\R}{\mathbb{R}}
\newcommand{\Q}{\mathbb{Q}}
\newcommand{\Z}{\mathbb{Z}}
\newcommand{\F}{\mathbb{F}}

\title{Homework 2}
\author{Asher Lucas-Cuddeback}

\begin{document}
\maketitle
\setcounter{problemcounter}{2}
\begin{problem}
  \begin{parts}

  \item Give a more explicit description of the group of square classes over
    the rationals, i.e.\ describe the group $\mathbb{Q}^*/(\mathbb{Q}^*)^2$.

    \begin{solution}
      The group, \( \Q^* \), is the multiplicative group of rational
      numbers without zero. Any element in this group is represented
      as, for \( a,b\in\Z \), \( \frac{a}{b} \). Via the fundamental
      theorem of arithmetic,
      \[
        \frac{a}{b}=(-1)^{e_{-1}}(2)^{e_2}(3)^{e_3}(5)^{e_5}(7)^{e_7}\dots\quad
        e_i\in\Z
      \]
      Since we care only about those factors which are numbers of the
      form \( e_i^{2n} \), \( n\in\Z \) and since for any number with
      an odd negative exponent can be written as \( e_i^{-2n-1} =
      e_i\cdot e_i^{-2n} \) the group of square classes over the
      rationals is the numbers described by:
      \[
        \Q^* / (\Q^*)^2=
        \left\{(-1)^{e_{-1}}(2)^{e_2}(3)^{e_3}(5)^{e_5}(7)^{e_7}\dots\,|\,
        e_i\in\{0, 1\}\right\}
      \]
      In this group, every element is its own inverse since for \(
      x\in\Q^*/(\Q^*)^2 \), \( x\cdot x = x^2 = 1 \).
    \end{solution}
    \newpage
  \item Use part (a), with our theory of non-degenerate symmetric bilinear
    forms from class, to argue that any such form on $V = \mathbb{Q}^n$ is
    congruent to a certain reduced diagonal representative. Is it
    possible for a
    given form to have two different diagonal representatives?

    \begin{solution}
      Our theory on symmetric, non-degenerate bilinear forms says
      that for any such form on some vector space over some field \(
      F^n \), there exists a basis for which the matrix
      representation of that form is diagonal. This is a universal
      property for \emph{any} such vector space and so holds in the
      case of \( \Q^n \). Now, regarding two reduced diagonal forms
      representing the same bilinear form we assert that it is possible.
      Take the non-degenerate, symmetric bilinear form with matrix
      representation \(
        \begin{bmatrix}
          2 & -6 \\
          -6 & 3
      \end{bmatrix} \) we can reduce it as follows:
      \[
        \begin{bmatrix}
          2 & -6 \\
          -6 & 3
        \end{bmatrix}\xrightarrow{R_2 + 3R_1}
        \begin{bmatrix}
          2 & -6 \\
          0 & -15
        \end{bmatrix}\xrightarrow{C_2 + 3C_1}
        \begin{bmatrix}
          2 & 0 \\
          0 & -15
        \end{bmatrix}
      \]
      This matrix cannot be reduced any further since both elements
      are in the group of square classes over the rationals. This
      representation is also congruent to:
      \[
        \begin{bmatrix}
          2 & -6 \\
          -6 & 3
        \end{bmatrix}\xrightarrow{R_1 + 2R_2}
        \begin{bmatrix}
          -10 & 0 \\
          -6 & 3
        \end{bmatrix}\xrightarrow{C_1 + 2C_2}
        \begin{bmatrix}
          -10 & 0 \\
          0 & 3
        \end{bmatrix}
      \]
      which is similarly fully reduced. So yes, it is possible for
      one symmetric, non-degenerate bilinear to have more than one
      diagonal representative.
    \end{solution}

  \end{parts}
\end{problem}
\newpage
\begin{problem}[Congruence classes on $F^2$]
  Describe the possible congruence classes of the matrix representation of a
  non-degenerate, symmetric, bilinear form on $F^2$ in the various cases.
  \begin{parts}

  \item $F = \mathbb{R}$
    \begin{solution}
      The only unique reduced squares in \( \R \) are
      \[
        \{1,-1, 0\}
      \]
      Since the bilinear form is non-degenerate, we cannot have \( 0
      \) in the main diagonal. Also, since it's symmetric, there
      exists a basis in which the matrix representation of the
      bilinear form is diagonal. Thus we have only three options:
      \[
        B\in\left\{
          \begin{bmatrix}
            1 & 0 \\
            0 & 1
          \end{bmatrix},
          \begin{bmatrix}
            1 & 0 \\
            0 & -1
          \end{bmatrix},
          \begin{bmatrix}
            -1 & 0 \\
            0 & -1
          \end{bmatrix}
        \right\}
      \]
    \end{solution}
    \newpage
  \item $F = \mathbb{C}$
    \begin{solution}
      There is only a single reduced value in the complex numbers.
      This is because \( \sqrt{-1}\in\C \) thus there is exactly one
      congruence class:
      \[
        A=
        \begin{bmatrix}
          1 & 0 \\
          0 & 1
        \end{bmatrix}
      \]
    \end{solution}

  \item $F = \mathbb{F}_3$
    \begin{solution}
      The field \( \F_3=\{0, 1, 2\} \). Again, since the bilinear
      form is symmetric and non-degenerate we can only have non-zero
      values on the main diagonal and zeros everywhere else. So for
      \( A\in\F_3^2 \),
      \[
        A\in\left\{
          \begin{bmatrix}
            1 & 0 \\
            0 & 1
          \end{bmatrix},
          \begin{bmatrix}
            1 & 0 \\
            0 & 2
          \end{bmatrix},
          \begin{bmatrix}
            2 & 0 \\
            0 & 2
          \end{bmatrix}
        \right\}
      \]
      Which we might assume completely characterizes the form.
      However, using \( B=
        \begin{bmatrix}1 & -1 \\ 1 & 1
      \end{bmatrix}\) we see that:
      \[
        \begin{bmatrix}
          2 & 0 \\
          0 & 2
        \end{bmatrix} =
        \begin{bmatrix}
          1 & 1 \\ -1 & 1
        \end{bmatrix}
        \begin{bmatrix}
          1 & 0 \\
          0 & 1
        \end{bmatrix}
        \begin{bmatrix}
          1 & -1 \\
          1 & 1
        \end{bmatrix}
      \]
      Which means the two are actually the same equivalence class. So
      the two classes are actually
      \[
        A\in\left\{
          \begin{bmatrix}
            1 & 0 \\
            0 & 1
          \end{bmatrix},
          \begin{bmatrix}
            1 & 0 \\
            0 & 2
          \end{bmatrix},
        \right\}
      \]
    \end{solution}
    \newpage
  \item The form is skew-symmetric and $F$ is any of the fields above.
    \begin{solution}
      There is only one reduced, non-degenerate, skew-symmetric
      bilinear form when \( F^n=F^{2k} \) is an even dimensional vector
      space. This form is represented in \( F^2 \) by:
      \[
        \begin{bmatrix}
          0 & -1
          1 & 0
        \end{bmatrix}
      \]
      So for every field discussed, this is the only form.
    \end{solution}

  \end{parts}
\end{problem}
\newpage
\begin{problem}
  \begin{parts}

  \item Consider the following bilinear form on $\mathbb{R}^2$:
    \[
      B\!\left(
        \begin{bmatrix}a\\b
        \end{bmatrix},
        \begin{bmatrix}c\\d
      \end{bmatrix}\right)
      = ac - ad - bc + 2bd.
    \]
    Is this form congruent to the usual dot product? If yes, find a new basis
    where this form is represented by the usual dot product.

    \begin{solution}
      The matrix representation of \( B \) is:
      \[
        B\left(
          \begin{bmatrix}
            a \\ b
          \end{bmatrix},
          \begin{bmatrix}
            c \\ d
        \end{bmatrix}\right) =
        \begin{bmatrix}
          a & b
        \end{bmatrix}
        \begin{bmatrix}
          1 & -1 \\
          -1 & 2
        \end{bmatrix}
        \begin{bmatrix}
          c \\ d
        \end{bmatrix}
      \]
      So if we decompose this with \( R_2 + R_1 \) and then \( C_2 +
      C_1 \) we get:
      \[
        \begin{bmatrix}
          1 & -1 \\
          -1 & 2
        \end{bmatrix}\sim
        \begin{bmatrix}
          1 & 0 \\
          0 & 1
        \end{bmatrix}
      \]
      The equation which describes these moves is:
      \[
        \begin{bmatrix}
          1 & 1 \\
          0 & 1
        \end{bmatrix}^T
        \begin{bmatrix}
          1 & -1 \\
          -1 & 2
        \end{bmatrix}
        \begin{bmatrix}
          1 & 1 \\
          0 & 1
        \end{bmatrix}=
        \begin{bmatrix}
          1 & 0 \\
          0 & 1
        \end{bmatrix}
      \]
      So by reading off the columns of the non-transposed matrix we
      get that the basis in which \( B \) is the normal dot product is
      \[
        \left\{
          \begin{bmatrix}
            1 \\ 0
          \end{bmatrix},
          \begin{bmatrix}
            1 \\ 1
          \end{bmatrix}
        \right\}
      \]
    \end{solution}
    \newpage

  \item Consider the following bilinear form on $\mathbb{R}^2$:
    \[
      B\!\left(
        \begin{bmatrix}a\\b
        \end{bmatrix},
        \begin{bmatrix}c\\d
      \end{bmatrix}\right)
      = 2ac + 3ad + 3bc + bd.
    \]
    Is this form congruent to the usual dot product? If yes, find a new basis
    where this form is represented by the usual dot product.

    \begin{solution}
      This form is represented by the matrix:
      \[
        B=
        \begin{bmatrix}
          2 & 3 \\
          3 & 1
        \end{bmatrix}
      \]
      Reducing this yields
      \[
        \begin{bmatrix}
          2 & 3 \\
          3 & 1
        \end{bmatrix}\xrightarrow{R_1 - 3R_2}
        \begin{bmatrix}
          -7 & 0 \\
          3 & 1
        \end{bmatrix}\xrightarrow{C_1 - 3C_2}
        \begin{bmatrix}
          -7 & 0 \\
          0 & 1
        \end{bmatrix}\xrightarrow[\frac{1}{\sqrt{7}}C_1]{\frac{1}{\sqrt{7}}R_1}
        \begin{bmatrix}
          -1 & 0 \\
          0 & 1
        \end{bmatrix}
      \]
      And since \( \sqrt{-1}\not\in\R \), we cannot reduce further.
      Thus, \( B \) is \emph{not} congruent to the normal dot product.
    \end{solution}

  \end{parts}
\end{problem}

\end{document}
